% golden B — 人手序列化（ir2tikz 的对照真值）
% source: tests/golden/B/ir.json  (single-stage amp, PLAN.md §6.6 金样)
% nets:
%   N_IN: port a(0,4), R1.1
%   N_B : R1.2, C1.1, Q1.B          <- junction (3,4)
%   N_C : Q1.C, R2.1
%   VCC : R2.2, VCC1.p
%   GND : C1.2, Q1.E, GND1.p, GND2.p
\documentclass[margin=5pt]{standalone}
\usepackage[american]{circuitikz}
\ctikzset{bipoles/length=1.0cm}
\begin{document}
\begin{circuitikz}

%% ==== [region] input-coupling ====
\coordinate (N_B) at (3,4);
\draw (1,4) to[R, l=$R_1$, a=$10\mathrm{k}\Omega$, name=R1] (3,4);
\draw (3,4) to[C, l=$C_1$, a=$100\mathrm{nF}$, name=C1] (3,2);
\node[ground] at (3,2) {};

%% ==== [region] amplifier ====
\node[npn] (Q1) at (5,4) {};
\node[anchor=west] at (5.2,4) {$Q_1$};
\draw (5,6) to[R, l_=$R_C$, a^=$2\mathrm{k}\Omega$, name=R2] (5,7.5);
\node[vcc] at (5,7.5) {};
\node[anchor=south] at (5,8) {$+12\mathrm{V}$};
\node[ground] at (5,2) {};

%% ==== wires ====
% net N_IN
\draw (0,4) -- (1,4);   % W1: port a -> R1.1
% net N_B
\draw (N_B) -- (Q1.B);  % W2: R1.2 -> Q1.B
% net N_C
\draw (Q1.C) -- (5,6);  % W3: Q1.C -> R2.1
% net GND
\draw (Q1.E) -- (5,2);  % W4: Q1.E -> GND2

%% ==== junctions / terminals / texts ====
\node[circ] at (N_B) {};
\node[ocirc] at (0,4) {};
\node[anchor=east] at (-0.3,4) {$a$};
\node[anchor=south east] at (0.2,4.6) {$u_i$};

\end{circuitikz}
\end{document}
